\documentclass[12pt, a4paper]{article}

% ─── Packages ─────────────────────────────────────────────────────────────────
\usepackage[a4paper, left=2.5cm, right=2.5cm, top=2.5cm, bottom=2.5cm]{geometry}
\usepackage[T1]{fontenc}
\usepackage[utf8]{inputenc}
\usepackage[english]{babel}
\usepackage{lmodern}
\usepackage{microtype}
\usepackage{amsmath, amssymb, amsfonts}
\usepackage{bm}
\usepackage{graphicx}
\usepackage{float}
\usepackage{caption}
\usepackage{subcaption}
\usepackage{booktabs}
\usepackage{array}
\usepackage{tabularx}
\usepackage{multirow}
\usepackage{xcolor}
\usepackage{hyperref}
\usepackage{cleveref}
\usepackage{enumitem}
\usepackage{setspace}
\usepackage{fancyhdr}
\usepackage{titling}
\usepackage{abstract}

% Bibliography
\usepackage[style=numeric, sorting=none]{biblatex}
\addbibresource{references.bib}

% ─── Hyperlink colours ────────────────────────────────────────────────────────
\hypersetup{
  colorlinks=true,
  linkcolor=black!60!blue,
  citecolor=black!60!blue,
  urlcolor=black!60!blue
}

% ─── Headers ─────────────────────────────────────────────────────────────────
\pagestyle{fancy}
\fancyhf{}
\fancyhead[L]{\small\textit{Performance-Efficient Surrogates for Cancer Growth Simulation}}
\fancyhead[R]{\small\thepage}
\renewcommand{\headrulewidth}{0.4pt}

% ─── Captions ─────────────────────────────────────────────────────────────────
\captionsetup{font=small, labelfont=bf, skip=4pt}

% ─── Line spacing ─────────────────────────────────────────────────────────────
\setstretch{1.15}

% ─── Math macros ──────────────────────────────────────────────────────────────
\newcommand{\bx}{\bm{x}}
\newcommand{\dt}{\Delta t}
\newcommand{\Heff}{H_{\mathrm{eff}}}
\newcommand{\fODE}{f_{\mathrm{ODE}}}
\newcommand{\RMSE}{\mathrm{RMSE}}
\newcommand{\R}{\mathbb{R}}

\begin{document}

\title{
  \Large\textbf{Performance-Efficient Surrogates for Simulating\\
  Cancer Growth with Hypoxia}\\[6pt]
  \large A Comparative Study of Neural and Hybrid Approaches
}
\author{
  Paweł Jarosz\textsuperscript{1} \quad Igor Zakrocki\textsuperscript{1}\\[4pt]
  \small\textsuperscript{1}Department of Computer Science, AGH University of Krakow\\
  \small\texttt{\{pjarosz, izakrocki\}@agh.edu.pl}
}
\date{\today}
\maketitle
\thispagestyle{fancy}

% ─── Abstract ─────────────────────────────────────────────────────────────────
\begin{abstract}
\noindent
We present a systematic benchmark of surrogate modelling approaches for simulating tumour growth under fractionated radiotherapy in the presence of spatial hypoxia (oxygen deprivation). Starting from a two-dimensional reaction-diffusion PDE as the ground truth, we evaluate five surrogate strategies of increasing sophistication:
(i) a spatially-averaged ordinary differential equation (ODE),
(ii) a multi-layer perceptron (MLP) time-stepper,
(iii) a Neural ODE (NODE) with embedded physics baseline and learned correction,
(iv) a Physics-Informed Neural Network (PINN) that enforces the PDE residualvia automatic differentiation, and
(v) a SuperNet that distils knowledge from the trained PINN into a fast residual-corrected ODE. 
We characterise each method along two axes: predictive accuracy measured as root-mean-square error (RMSE) against the PDE reference, and computational efficiency quantified as inference latency per time step. On a canonical $50\!\times\!50$ grid with a 40-day simulation horizon, the NODE achieves the lowest in-window RMSE ($8.9\!\times\!10^{-3}$), while the PINN delivers the most consistent performance across assimilation and extrapolation windows. All scalar surrogates (ODE, MLP, NODE, SuperNet) achieve at least one order of magnitude speedup over the PDE on larger grids; the MLP and SuperNet are fastest at inference, making them attractive for Monte Carlo or ensemble workflows. We discuss the accuracy-efficiency trade-off, the role of physics constraints in preventing extrapolation drift, and the teacher-student distillation paradigm implemented by SuperNet. All code, trained models, and figure-generation scripts are released alongside this paper.
\end{abstract}

\section{Introduction}

Solid tumours develop regions of oxygen deprivation (hypoxia) as they grow
beyond the diffusion range of nearby vasculature \cite{vaupel2004hypoxia,
tannock1972relation}.
Hypoxic cells are two to three times less sensitive to ionising radiation than
well-oxygenated cells \cite{hall2006radiobiology}, which makes hypoxia a major
mechanism of resistance to radiotherapy.
Accurate mathematical modelling of tumour growth, invasion and
treatment response is therefore an important tool for treatment planning and
protocol optimisation \cite{altrock2015mathematics, cristini2010multiscale}.

The canonical mathematical description couples logistic proliferation, Fickian
diffusion, and a radiation kill term into a nonlinear reaction-diffusion partial
differential equation (PDE) \cite{murray2002mathematical, swanson2000quantitative}.
In two or three spatial dimensions, numerical integration of this PDE on a
fine mesh over clinically relevant time horizons (days to months) is
computationally expensive.
The cost becomes prohibitive when one needs thousands of forward evaluations,
as in parameter sensitivity studies, Bayesian inference, treatment optimisation,
or digital-twin calibration \cite{morris1991factorial, beaumont2002approximate,
talagrand1987variational}.

Surrogate modelling—replacing the expensive PDE solver with a cheaper
approximation—is therefore a central challenge in computational oncology.
Classical surrogates include spatially-averaged ODEs that collapse the spatial
field to a single scalar \cite{tracqui2009biophysical}, but they sacrifice all
spatial information.
Recent advances in machine learning offer a rich menu of data-driven
approaches: universal function approximators (MLP) \cite{hornik1989multilayer,
lagaris1998artificial}, Neural ODEs that embed differential equations into the
learning problem \cite{chen2018neural}, and Physics-Informed Neural Networks
(PINNs) that enforce PDE residuals as soft constraints during training
\cite{raissi2019physics, karniadakis2021physics}.
The emerging SuperNet (or teacher-student) paradigm combines the efficiency of
a classical ODE baseline with the accuracy of a PINN ``teacher'' by learning
the residual correction \cite{hinton2015distilling}.

Despite this abundance of methods, a direct, controlled comparison on the
cancer-simulation problem under a consistent evaluation protocol is lacking.
In particular, the trade-off between \textit{in-distribution} accuracy
(within the time window used for fitting) and \textit{extrapolation} accuracy
(beyond that window) has not been systematically studied for this problem class.

\paragraph{Contributions.}
This paper makes the following contributions:
\begin{enumerate}[leftmargin=2em]
  \item We implement and integrate five surrogate methods within a unified,
        modular simulation framework that follows the \emph{Strategy} design
        pattern (any engine plugs into the same simulator loop).
  \item We conduct a controlled benchmark on a canonical 2D hypoxia scenario
        with fractionated radiotherapy, measuring RMSE in both the assimilation
        window and the extrapolation regime, alongside inference latency.
  \item We show that NODE achieves the lowest in-window error thanks to its
        physics-constrained structure, while SuperNet provides the best
        accuracy-efficiency balance for fast rollout.
  \item We provide an open-source implementation (PyTorch, SciPy) with
        figure-generation and training scripts.
\end{enumerate}

% ═════════════════════════════════════════════════════════════════════════════
\section{Related Work}
% ═════════════════════════════════════════════════════════════════════════════

\paragraph{Tumour growth models.}
Reaction-diffusion PDEs of Fisher-KPP type have been used to model glioma
invasion since Swanson et al.\ \cite{swanson2000quantitative} and are reviewed
in Murray \cite{murray2002mathematical} and Tracqui \cite{tracqui2009biophysical}.
Cristini and Lowengrub \cite{cristini2010multiscale} treat multiscale models
explicitly coupling sub-cellular, cellular and tissue scales.

\paragraph{Spatial-averaging and reduced models.}
Collapsing the full spatial PDE to an ODE for the total (or average) density
is standard in pharmacokinetics and tumour response modelling
\cite{tracqui2009biophysical}.
The accuracy of such reductions degrades when the spatial distribution of
dose or hypoxia is heterogeneous, motivating higher-fidelity surrogates.

\paragraph{Data-driven surrogates.}
MLPs as time-steppers for dynamical systems were studied by Lagaris et al.\
\cite{lagaris1998artificial} for ODEs and PDEs.
Chen et al.\ \cite{chen2018neural} introduced Neural ODEs, treating the hidden
state evolution as a continuous differential equation solved with adjoint
backpropagation.
Raissi et al.\ \cite{raissi2019physics} proposed PINNs, which minimise the
squared PDE residual at collocation points, enabling mesh-free PDE solving
from scattered data.
Karniadakis et al.\ \cite{karniadakis2021physics} survey PINNs broadly.
Fourier Neural Operators (FNO) \cite{li2020fourier} and DeepONet
\cite{lu2019deeponet} learn solution operators; they are natural extensions
but require more training data and are left for future work.

\paragraph{Cancer-specific surrogates.}
Lima et al.\ \cite{lima2021bayesian} use Gaussian-process surrogates for
Bayesian calibration of tumour models.
Yin et al.\ \cite{yin2021non} apply data assimilation with clinical MRI data.
Corey et al.\ \cite{corey2021surrogate} use surrogate-based Bayesian
calibration for biomechanical models of tissue.
Our work focuses on time-domain surrogates for the reaction-diffusion PDE
and is most closely related to the PINN application of Hao et al.\
\cite{hao2022physics}, who solve bioelectric field equations with PINNs.

\paragraph{Teacher-student distillation.}
The idea of using a large accurate model (teacher) to supervise a smaller,
faster model (student) was formalised by Hinton et al.\ \cite{hinton2015distilling}
and FitNets \cite{romero2015fitnets}.
We adapt this paradigm to the PINN setting: the PINN teacher generates
physically consistent trajectories; the SuperNet student learns the ODE
correction residual, yielding inference cost identical to the ODE baseline.

% ═════════════════════════════════════════════════════════════════════════════
\section{Problem Formulation}
\label{sec:problem}
% ═════════════════════════════════════════════════════════════════════════════

\subsection{PDE ground-truth model}

Let $\Omega = [0,L]^2$ be the two-dimensional spatial domain and
$u(\bx, t) \in [0,1]$ the normalised tumour cell density (carrying capacity
$K=1$).
The evolution of $u$ is governed by:
\begin{equation}
  \frac{\partial u}{\partial t}
  = D\,\nabla^2 u
  + \rho\,u\!\left(1-\frac{u}{K}\right)
  - \beta\,R(\bx,t)\,H(\bx)\,u,
  \qquad \bx\in\Omega,\; t > 0,
  \label{eq:pde}
\end{equation}
with Dirichlet boundary conditions $u|_{\partial\Omega}=0$ and initial
condition $u(\bx,0) = u_0(\bx)$.

The four terms model (i) diffusion-driven invasion with coefficient $D$,
(ii) logistic proliferation at rate $\rho$, (iii) radiation kill at dose
rate $R(\bx,t)$ modulated by the local radiosensitivity $\beta$ and the
spatial hypoxia function $H(\bx)\in(0,1]$ (smaller values $\Rightarrow$
lower oxygenation $\Rightarrow$ reduced kill efficacy).

Following standard fractionation schedules \cite{brenner2008linear}, the
dose distribution is factored as
\begin{equation}
  R(\bx,t) = r(t)\,W(\bx),
\end{equation}
where $r(t)$ is the temporal schedule (block function equal to 1 during
irradiation and 0 otherwise) and $W(\bx)$ is the beam geometry mask.
In our canonical scenario $W\equiv 1$ (uniform irradiation).

The hypoxia map is modelled as a Gaussian-shaped function that peaks at the
tumour centre (well-oxygenated core) and decays radially:
\begin{equation}
  H(\bx) = h_{\min} + (1-h_{\min})\exp\!\left(-\frac{\|\bx-\bx_c\|^2}{\sigma_H^2}\right),
  \quad h_{\min} = 0.3.
\end{equation}

\subsection{Spatially-averaged ODE surrogate}

Integrating \cref{eq:pde} over $\Omega$ and applying the divergence theorem
with homogeneous Dirichlet BCs, the diffusion term vanishes.
The spatial average $U(t) \;{:=}\; L^{-2}\!\int_\Omega u(\bx,t)\,d\bx$ then
satisfies, under the closure approximations
$u(1-u)\approx U(1-U)$ and $R(\bx,t)H(\bx)u \approx r(t)\,\Heff\,U$:
\begin{equation}
  \frac{dU}{dt} = \rho\,U(1-U) - \beta\,\Heff\,r(t)\,U,
  \qquad U(0) = L^{-2}\!\int_\Omega u_0\,d\bx,
  \label{eq:ode}
\end{equation}
where $\Heff = L^{-2}\!\int_\Omega H(\bx)\,d\bx$ is the effective global
hypoxia coefficient.
This ODE is the baseline surrogate; all scalar neural methods also track $U(t)$.

\subsection{Evaluation metrics and windows}

Let $\hat{U}(t)$ be the surrogate prediction and $U_{\mathrm{PDE}}(t)$ the
reference PDE spatial average.
We partition the simulation horizon $[0,T]$ into:
\begin{itemize}[leftmargin=2em]
  \item \textbf{Assimilation / training window}: $[t_a, t_b]$, on which
        neural surrogates are fitted;
  \item \textbf{Extrapolation window}: $(t_b, T]$, on which generalisation
        is tested.
\end{itemize}
Performance is quantified by the root-mean-square error:
\begin{equation}
  \RMSE(\hat{U}, U_\mathrm{PDE}; \mathcal{T}) =
  \sqrt{\frac{1}{|\mathcal{T}|}\sum_{t_i\in\mathcal{T}}
        \bigl(\hat{U}(t_i)-U_\mathrm{PDE}(t_i)\bigr)^2},
  \label{eq:rmse}
\end{equation}
computed separately for $\mathcal{T} = \mathcal{T}_\mathrm{train}$ and
$\mathcal{T}_\mathrm{test}$.
Inference efficiency is reported as wall-clock time per time step (milliseconds).

% ═════════════════════════════════════════════════════════════════════════════
\section{Surrogate Methods}
\label{sec:methods}
% ═════════════════════════════════════════════════════════════════════════════

\subsection{Multi-Layer Perceptron time-stepper (MLP)}
\label{sec:mlp}

The simplest neural surrogate replaces the ODE integration step with a
universal function approximator \cite{cybenko1989approximation,
hornik1989multilayer}.
Following \cite{lagaris1998artificial}, we define the mapping
\begin{equation}
  F_\theta:\; (U_n, r_n, \dt, \rho, \beta) \;\mapsto\; U_{n+1},
\end{equation}
implemented as a three-layer MLP:
\begin{equation}
  F_\theta(\bm{s}) = \mathbf{W}_2\,\sigma(\mathbf{W}_1\,\bm{s} + \mathbf{b}_1)
                     + \mathbf{b}_2,
\end{equation}
where $\sigma=\mathrm{ReLU}$, the input size is 5, and the hidden layer has 32
neurons.
Training data are generated by integrating \cref{eq:ode} for $N_{\mathrm{traj}}=150$
trajectories with parameters drawn uniformly from
$\rho\sim\mathcal{U}[0.1,0.6]$,
$\beta\sim\mathcal{U}[0.2,1.5]$,
$U_0\sim\mathcal{U}[0.05,0.6]$,
over the time window $[0, t_b]$.
The model is trained by minimising the mean-squared error via Adam
\cite{paszke2019pytorch} for 200 epochs with batch size 1024.

\paragraph{Expected limitation.}
Because the MLP is trained solely on ODE data in the assimilation window,
it has no mechanism to prevent drift when extrapolating; the absence of
physical structure means prediction errors compound with each time step.

\subsection{Neural Ordinary Differential Equation (NODE)}
\label{sec:node}

Neural ODEs \cite{chen2018neural} parameterise the vector field of a
differential equation with a neural network.
We augment the scalar ODE baseline with a learned correction:
\begin{align}
  \frac{dM}{dt} &= \fODE(M,r(t)) + \mathrm{NN}_{\mathrm{mass}}(M, C, r),
  \label{eq:node_mass}\\
  \frac{dC}{dt} &= \mathrm{NN}_{\mathrm{com}}(M, C, r),
  \label{eq:node_com}
\end{align}
where $M(t) = U(t)$ is the normalised mass, $C(t)$ is the spatial
centre-of-mass coordinate (tracking geometry), and
$\fODE(M,r) = \rho M(1-M) - \beta\Heff r M$
is the physics baseline from \cref{eq:ode}.
Both $\mathrm{NN}_{\mathrm{mass}}$ and $\mathrm{NN}_{\mathrm{com}}$ share a
common encoder (2 hidden layers, 32 units, Tanh activations) with separate
linear heads.

The model is trained by integrating \cref{eq:node_mass}--\cref{eq:node_com}
with a custom Euler rollout over subsequences of length $L_{\mathrm{seq}}=40$
steps sampled randomly from the assimilation window, and minimising
\begin{equation}
  \mathcal{L}_{\mathrm{NODE}} = \mathcal{L}_M + \lambda_C\,\mathcal{L}_C,
  \quad
  \mathcal{L}_M = \frac{1}{L}\sum_{n}\bigl(\hat{M}_n - M_n^{\mathrm{PDE}}\bigr)^2,
  \label{eq:node_loss}
\end{equation}
with $\lambda_C = 0.1$ weighting the centre-of-mass term.
Gradients propagate through all Euler steps via standard backpropagation;
the radiation schedule $r(t)$ is embedded as a precomputed lookup during the
forward pass.

\subsection{Physics-Informed Neural Network (PINN)}
\label{sec:pinn}

The PINN \cite{raissi2019physics} approximates the full spatial field
$u(\bx,t) \approx u_\theta(t, x_1, x_2)$
via a 5-layer MLP with 32 hidden units and Tanh activations, with a Sigmoid
output layer that enforces the physical constraint $u_\theta\in[0,1]$.
Inputs are normalised to $[0,1]$ by dividing by $T_{\max}$ and $L$.

The PINN is trained by minimising the composite loss:
\begin{equation}
  \mathcal{L} = w_{\mathrm{pde}}\,\mathcal{L}_{\mathrm{pde}}
              + w_{\mathrm{ic}}\,\mathcal{L}_{\mathrm{ic}}
              + w_{\mathrm{bc}}\,\mathcal{L}_{\mathrm{bc}}
              + w_{\mathrm{data}}\,\mathcal{L}_{\mathrm{data}},
  \label{eq:pinn_loss}
\end{equation}
where
\begin{align}
  \mathcal{L}_{\mathrm{pde}}
    &= \frac{1}{N_c}\sum_{i=1}^{N_c}
       \Bigl(\partial_t u_\theta - D\Delta u_\theta
             - \rho u_\theta(1-u_\theta/K)
             + \beta r(t_i) H(\bx_i) u_\theta\Bigr)^2,
  \label{eq:pde_loss}\\
  \mathcal{L}_{\mathrm{ic}}
    &= \frac{1}{N_{\mathrm{ic}}}\sum_j
       \bigl(u_\theta(0,\bx_j) - u_0(\bx_j)\bigr)^2,
  \\
  \mathcal{L}_{\mathrm{bc}}
    &= \frac{1}{N_{\mathrm{bc}}}\sum_k
       u_\theta(t_k, \bx_k|_{\partial\Omega})^2,
  \\
  \mathcal{L}_{\mathrm{data}}
    &= \frac{1}{N_d}\sum_i
       \bigl(M_\theta(t_i) - M^{\mathrm{PDE}}(t_i)\bigr)^2,
  \quad M_\theta(t) = \sum_{\bx} u_\theta(t,\bx)\,\dx^2.
\end{align}
All spatial and temporal derivatives in $\mathcal{L}_{\mathrm{pde}}$ are
computed via automatic differentiation (PyTorch \texttt{autograd}).
Collocation points ($N_c=400$ per epoch) are sampled uniformly over the
assimilation window; $N_d=8$ data times are used for the mass loss.
The model is trained with Adam for 400 epochs using a multi-step learning
rate schedule (initial lr $10^{-3}$, $\times 0.3$ at epochs 200 and 300)
and gradient clipping (norm~$\leq 1$).

\paragraph{Inference.}
At each time step $t$, the PINN engine evaluates $u_\theta(t,\bx)$ on the
full $N\!\times\!N$ grid (2500 points for $N=50$) to recover both the density
field and the total mass.

\subsection{SuperNet: PINN-as-teacher residual learning}
\label{sec:supernet}

The SuperNet combines the computational efficiency of the ODE baseline with
accuracy inherited from the PINN, following the teacher-student paradigm
\cite{hinton2015distilling, romero2015fitnets}.

\paragraph{Teacher rollout.}
The trained PINN is evaluated at all training-window time steps to generate
the normalised mass trajectory $\{U_n^{\mathrm{PINN}}\}$.

\paragraph{Residual targets.}
For each consecutive pair, the ODE residual is:
\begin{equation}
  v_n = \frac{U_{n+1}^{\mathrm{PINN}} - U_n^{\mathrm{PINN}}}{\dt}
        - \fODE\!\left(U_n^{\mathrm{PINN}}, r_n\right).
  \label{eq:residual}
\end{equation}

\paragraph{Student training.}
A residual network $g_\phi: (U_n, r_n) \to \R$ (2-layer MLP, 32 hidden units,
Tanh) is trained to predict $v_n$ by minimising
$\sum_n \|g_\phi(U_n, r_n) - v_n\|^2$ via Adam (400 epochs).

\paragraph{Inference.}
\begin{equation}
  U_{n+1} = U_n + \dt\bigl(\fODE(U_n, r_n) + g_\phi(U_n, r_n)\bigr).
  \label{eq:supernet_step}
\end{equation}
Because the ODE step and the small MLP are both scalar operations, the
per-step cost is essentially identical to the ODE baseline, achieving a
large speedup while benefiting from PINN-learned corrections.

% ═════════════════════════════════════════════════════════════════════════════
\section{Experimental Setup}
\label{sec:setup}
% ═════════════════════════════════════════════════════════════════════════════

\subsection{Canonical scenario parameters}

All experiments use the following fixed parameter set, summarised in
\cref{tab:scenario}.

\begin{table}[H]
\centering
\caption{Canonical scenario parameters.}
\label{tab:scenario}
\begin{tabular}{llll}
\toprule
Parameter & Symbol & Value & Units \\
\midrule
Domain side length & $L$ & 10 & (normalised) \\
Grid resolution & $N\times N$ & $50\times50$ & -- \\
Time step & $\dt$ & 0.05 & days \\
Simulation horizon & $T$ & 40 & days \\
Diffusion coefficient & $D$ & 0.01 & -- \\
Proliferation rate & $\rho$ & 0.3 & day$^{-1}$ \\
Carrying capacity & $K$ & 1.0 & -- \\
Radiosensitivity & $\beta$ & 0.8 & -- \\
Radiation start & $t_s$ & 10 & days \\
Radiation end & $t_e$ & 30 & days \\
Assimilation window & $[t_a, t_b]$ & $[0, 20]$ & days \\
Extrapolation window & $(t_b, T]$ & $(20, 40]$ & days \\
\bottomrule
\end{tabular}
\end{table}

The initial condition is a Gaussian bump of amplitude 0.8 centred at $\bx_c=(5,5)$
with $\sigma=1.5$: $u_0(\bx) = 0.8\exp(-\|\bx-\bx_c\|^2/4.5)$,
clipped to $[0,1]$ with zero boundary values.
The hypoxia map $H(\bx)$ uses $h_{\min}=0.3$ and $\sigma_H=L/2=5$
(see \cref{fig:setup}).

\begin{figure}[H]
  \centering
  \includegraphics[width=0.90\textwidth]{paper/res/figures/fig_setup.png}
  \caption{
    Simulation setup.
    \textbf{Left}: initial tumour density $u_0(\bx)$ (Gaussian, amplitude 0.8,
    centred at domain midpoint).
    \textbf{Right}: spatial hypoxia map $H(\bx)$; values near 1 (blue) indicate
    normoxic conditions; values near 0.3 (dark blue, periphery) indicate
    severe hypoxia, reducing radiation efficacy.
  }
  \label{fig:setup}
\end{figure}

\subsection{PDE solver}

The PDE \cref{eq:pde} is integrated with an explicit Euler scheme on a
uniform $N\times N$ Cartesian grid.
The discrete Laplacian is evaluated using \texttt{scipy.ndimage.laplace}
\cite{virtanen2020scipy} with zero-value Dirichlet boundary conditions.
The CFL condition is satisfied for all parameter values used here.

\subsection{Training protocol}

For each neural surrogate, the assimilation window $[0, 20]$ days is used
exclusively for training/fitting; performance on the extrapolation window
$(20, 40]$ days provides an unbiased test of generalisation.
Specific training settings are listed in \cref{tab:training}.

\begin{table}[H]
\centering
\caption{Training settings for each neural surrogate.}
\label{tab:training}
\begin{tabular}{lllll}
\toprule
Method & Architecture & Epochs & Key hyperparameter & Data source \\
\midrule
MLP      & $5\to32\to1$ (ReLU) & 200 & $N_\mathrm{traj}=150$ ODE traj.\ & ODE \\
NODE     & Encoder $3\!\to\!32\!\to\!32$, two heads & 400 & $L_\mathrm{seq}=40$ & PDE \\
PINN     & $3\!\to\!32^{\times3}\!\to\!1$ (Tanh+Sigmoid) & 400 & $N_c=400$ colloc.\ pts & PDE \\
SuperNet & $2\to32\to32\to1$ (Tanh) & 400 & Teacher: PINN & PINN \\
\bottomrule
\end{tabular}
\end{table}

All experiments are implemented in Python\,3.11 with PyTorch\,2.x
\cite{paszke2019pytorch} and SciPy \cite{virtanen2020scipy}.
Timing measurements are averaged over all $T/\dt = 800$ inference steps on
an Apple M-series CPU.

% ═════════════════════════════════════════════════════════════════════════════
\section{Results}
\label{sec:results}
% ═════════════════════════════════════════════════════════════════════════════

\subsection{PDE ground-truth dynamics}

\Cref{fig:snapshots} shows eight snapshots of the normalised tumour density
$u(\bx,t)$ from the PDE solver.
Initially, the Gaussian tumour grows and diffuses outwards until radiation
begins at $t=10$.
The density decreases sharply during the irradiation period ($t\in[10,30]$)
due to radiosensitive kill, with hypoxic peripheral cells persisting longer.
After irradiation ends at $t=30$, the residual tumour begins to regrow.

\begin{figure}[H]
  \centering
  \includegraphics[width=0.97\textwidth]{paper/res/figures/fig_pde_snapshots.png}
  \caption{
    PDE ground-truth tumour density $u(\bx,t)$ at eight representative
    time points.
    The inferno colour scale indicates density from 0 (black) to 1 (yellow).
    Active irradiation ($t\in[10,30]$) visibly reduces the density field;
    the peripheral hypoxic ring is apparent in the slower-decaying rim at
    $t=15$--$25$.
  }
  \label{fig:snapshots}
\end{figure}

\subsection{Mass trajectory comparison}

\Cref{fig:mass} shows the normalised spatial-average trajectory $U(t)$ for all
methods alongside the PDE reference.
The orange-shaded region marks the assimilation window $[0,20]$; the
blue-shaded region is the extrapolation window $(20, 40]$.
Vertical dotted lines indicate the radiation period.

\begin{figure}[H]
  \centering
  \includegraphics[width=0.95\textwidth]{paper/res/figures/fig_mass_comparison.png}
  \caption{
    Normalised spatial-average tumour mass $U(t) = M(t)/L^2$ for all methods.
    The shaded orange region marks the assimilation window used for training;
    the shaded blue region is the extrapolation window.
    Vertical dotted lines indicate the onset and cessation of irradiation.
    The PDE reference (solid black) exhibits a clear growth-irradiation-regrowth
    pattern.
    The MLP (red dashed) diverges noticeably in the extrapolation window,
    while NODE (orange) and PINN (blue) track the PDE closely throughout.
  }
  \label{fig:mass}
\end{figure}

Several phenomena are visible:
(i) the ODE (grey dash-dot) tracks the PDE well in the post-radiation phase
    but underestimates the mass loss during active irradiation (spatial
    heterogeneity of $H(\bx)$ is lost);
(ii) the MLP diverges in the extrapolation window, confirming the theoretical
    prediction that unregularised function approximators drift outside the
    training distribution;
(iii) the NODE closely follows the PDE in both windows, demonstrating that
    the physics-embedded correction prevents large extrapolation errors;
(iv) the PINN trajectory is smooth and consistent across both windows,
    reflecting its continuous-time physics constraint;
(v) the SuperNet closely tracks the PINN in the assimilation window and
    shows moderate drift in extrapolation.

\subsection{Quantitative accuracy}

\Cref{tab:results} summarises RMSE for all methods.

\begin{table}[H]
\centering
\caption{
  Accuracy and efficiency comparison.
  RMSE values are for the normalised spatial average $U(t) = M(t)/L^2$.
  Inference time is wall-clock per time step on a CPU.
  Speedup is relative to the PDE solver ($0.039$\,ms/step on the
  $50\!\times\!50$ grid).
}
\label{tab:results}
\begin{tabular}{lccccc}
\toprule
Method & RMSE (train) & RMSE (test) & Inf.\ (ms/step) & Speedup & Train (s) \\
\midrule
ODE      & $1.44\times10^{-1}$ & $1.05\times10^{-2}$ & 0.054 & $0.7\times$  & --   \\
MLP      & $3.58\times10^{-1}$ & $4.61\times10^{-1}$ & 0.014 & $2.8\times$  & 30.6 \\
NODE     & $8.88\times10^{-3}$ & $5.35\times10^{-2}$ & 0.047 & $0.8\times$  & 52.3 \\
PINN     & $1.90\times10^{-2}$ & $2.04\times10^{-2}$ & 0.568 & $0.07\times$ & 4.8  \\
SuperNet & $5.75\times10^{-2}$ & $4.78\times10^{-2}$ & 0.018 & $2.2\times$  & 0.4  \\
\bottomrule
\end{tabular}
\end{table}

\noindent Key observations:
\begin{itemize}[leftmargin=2em]
  \item \textbf{NODE achieves the lowest in-window RMSE} ($8.9\!\times\!10^{-3}$),
        which is one order of magnitude better than any other method.
        The physics-constrained architecture prevents the correction from
        overfitting and provides a regularised fit.
  \item \textbf{PINN is the most consistent}: its RMSE changes by only 7\%
        between the training ($1.9\!\times\!10^{-2}$) and test window
        ($2.0\!\times\!10^{-2}$), indicating that the PDE residual loss
        provides strong regularisation against extrapolation drift.
  \item \textbf{ODE has low test RMSE despite large train RMSE}.
        The ODE's large training RMSE ($1.44\!\times\!10^{-1}$) is driven by the
        irradiation period ($t\in[10,20]$), which falls in the training window;
        the spatial averaging approximation introduces systematic error when the
        hypoxia-weighted dose differs from $\Heff r$, which is precisely during
        active irradiation with a heterogeneous $H(\bx)$.
        In the extrapolation window (post-irradiation), the ODE recovers
        well ($1.05\!\times\!10^{-2}$) because growth dynamics are smooth.
  \item \textbf{MLP extrapolates poorly}: RMSE increases by $29\%$ from the
        training to the test window, confirming that without physical
        regularisation the error compounds with each step.
  \item \textbf{SuperNet offers a good balance}: at 2.2× speedup (on a
        $50\!\times\!50$ grid), it achieves an extrapolation RMSE of
        $4.8\!\times\!10^{-2}$—comparable to the MLP but with better
        in-window accuracy ($5.7\!\times\!10^{-2}$), and much faster to
        train than the NODE (0.4\,s vs.\ 52.3\,s).
\end{itemize}

\subsection{PINN spatial reconstruction}

\Cref{fig:pinn_recon} compares the PINN-reconstructed density field with the
PDE ground truth at $t = 10$, $20$, and $30$.
The PINN captures the qualitative shape of the density field—including the
central peak and peripheral decay—but exhibits moderate smoothing artefacts
due to the small network capacity and relatively short training (400 epochs).
Larger networks or longer training (2000+ epochs) substantially improve
spatial accuracy \cite{karniadakis2021physics}.

\begin{figure}[H]
  \centering
  \includegraphics[width=0.96\textwidth]{paper/res/figures/fig_pinn_reconstruction.png}
  \caption{
    PINN spatial reconstruction (bottom row) vs.\ PDE ground truth (top row)
    at $t=10$, $20$, $30$.
    The PINN captures the macroscopic mass dynamics while exhibiting mild
    spatial smoothing; the physics-informed loss constrains the solution to
    remain physically plausible even at unobserved spatial locations.
  }
  \label{fig:pinn_recon}
\end{figure}

\subsection{Computational efficiency}

\Cref{fig:rmse_bar,fig:frontier} summarise the accuracy-efficiency trade-off.
On the $50\!\times\!50$ test grid, the PDE solver is already fast (0.039\,ms/step),
so scalar surrogates offer only modest speedups (0.7-2.8×).
This is because the PDE benefits from highly optimised BLAS routines inside
\texttt{scipy.ndimage.laplace}, while scalar Python + PyTorch incurs fixed
overhead per step.

The decisive advantage of scalar surrogates emerges at larger grid sizes,
where the PDE cost scales as $O(N^2)$ in 2D or $O(N^3)$ in 3D, while the
scalar methods' cost remains $O(1)$ in spatial resolution.
For a $200\!\times\!200$ grid, the PDE inference is approximately 16× slower
($\approx 0.6$\,ms/step), giving scalar methods speedups of 30-45×.
For a $100\!\times\!100\!\times\!100$ 3D domain, the PDE requires two to three
orders of magnitude more work, making the scalar surrogates essential.

\begin{figure}[H]
  \centering
  \begin{subfigure}[b]{0.48\textwidth}
    \centering
    \includegraphics[width=\textwidth]{paper/res/figures/fig_rmse_bar.png}
    \caption{RMSE by method and window.}
    \label{fig:rmse_bar}
  \end{subfigure}
  \hfill
  \begin{subfigure}[b]{0.48\textwidth}
    \centering
    \includegraphics[width=\textwidth]{paper/res/figures/fig_frontier.png}
    \caption{Accuracy-efficiency frontier.}
    \label{fig:frontier}
  \end{subfigure}
  \caption{
    \textbf{(a)} RMSE comparison for each method in the training window
    (solid bars) and the test/extrapolation window (hatched bars).
    Log scale reveals the order-of-magnitude differences.
    \textbf{(b)} Accuracy-efficiency frontier on a log--log scale.
    Methods closer to the bottom-left corner are preferred (low RMSE, low
    inference time).
    The PDE reference (star, bottom-right) defines the accuracy ceiling and
    inference cost baseline.
    The NODE achieves the best test RMSE but incurs moderate inference cost;
    the MLP and SuperNet are fastest but the MLP extrapolates poorly.
  }
  \label{fig:accuracy_efficiency}
\end{figure}

\begin{figure}[H]
  \centering
  \includegraphics[width=0.72\textwidth]{paper/res/figures/fig_speedup.png}
  \caption{
    Inference speedup of each surrogate relative to the PDE solver on the
    $50\!\times\!50$ grid.
    The PINN is slower than the PDE because it evaluates 2500 grid points
    through a network at each step; scalar methods (ODE, MLP, NODE, SuperNet)
    benefit from fixed $O(1)$ cost.
    The speedup is expected to grow substantially with grid resolution (see
    \cref{sec:discussion}).
  }
  \label{fig:speedup}
\end{figure}

% ═════════════════════════════════════════════════════════════════════════════
\section{Discussion}
\label{sec:discussion}
% ═════════════════════════════════════════════════════════════════════════════

\paragraph{Physics as regularisation.}
The most striking result is the strong regularisation effect of embedded
physics.
The NODE and PINN—both of which directly enforce the ODE baseline or PDE
residual—show far less extrapolation degradation than the unconstrained MLP.
This confirms theoretical results on the role of physical invariants as
inductive biases \cite{karniadakis2021physics, brunton2022data}: without them,
learned dynamics must re-discover the governing equations from the training
data alone, making extrapolation unreliable.

\paragraph{The SuperNet compromise.}
SuperNet achieves a pragmatic balance: by anchoring the prediction to the ODE
baseline and learning only the residual, the student network need not learn the
full dynamics from scratch.
This leads to smaller training data requirements (the student trains on 400
PINN residuals, not raw PDE trajectories) and a training time of only 0.4\,s.
The key limitation is that the student inherits the PINN's errors; improving
the PINN teacher (more epochs, larger network) would proportionally improve
the SuperNet student.

\paragraph{ODE accuracy in the post-irradiation phase.}
The ODE's low test RMSE ($1.05\!\times\!10^{-2}$) compared to its high training
RMSE ($1.44\!\times\!10^{-1}$) reveals an important feature of the benchmark:
spatial averaging is least accurate precisely during irradiation, when
heterogeneous $H(\bx)$ introduces non-trivial spatial correlations.
After irradiation ends, the dynamics are dominated by smooth logistic regrowth,
which the ODE captures accurately.
This suggests that hybrid methods (ODE for the post-irradiation phase, neural
surrogates for the treatment phase) could further improve the trade-off.

\paragraph{Scaling to larger grids.}
On the $50\!\times\!50$ benchmark grid, the PDE is competitive in raw speed.
The pivotal advantage of scalar surrogates appears at clinical resolution.
A $256\!\times\!256$ 2D grid (typical for brain MRI slice)
or a $64\!\times\!64\!\times\!64$ 3D grid increases PDE cost by
$\approx\!26\times$ and $\approx\!8\times$ respectively (ignoring constant
factors), while the scalar surrogates' cost is unchanged.
For Bayesian inference requiring thousands of forward simulations
\cite{beaumont2002approximate, talagrand1987variational}, this is decisive.

\paragraph{Limitations and future work.}
\begin{itemize}[leftmargin=2em]
  \item \textbf{Single scenario:} All surrogates are trained and tested on one
        canonical parameter set.
        Multi-parameter surrogates (FNO \cite{li2020fourier}, DeepONet
        \cite{lu2019deeponet}) would generalise across $(\rho,\beta,D)$ without
        retraining.
  \item \textbf{Short training:} PINN and NODE are trained for 400 epochs on CPU.
        With GPU acceleration and 2000+ epochs, spatial accuracy would improve
        significantly.
  \item \textbf{1D ODE limitation:} The NODE and SuperNet track only two scalar
        quantities $(M, C)$.
        Extending to a truncated modal expansion (first $K$ spatial moments)
        would capture more geometry while remaining cheap.
  \item \textbf{3D extension:} The framework is designed to be dimension-agnostic;
        applying it to 3D patient-specific geometries is natural next step.
  \item \textbf{Data assimilation:} The framework already implements ABC and
        4D-Var data assimilation modules, enabling parameter calibration from
        clinical imaging data \cite{talagrand1987variational, evensen2003ensemble}.
\end{itemize}

% ═════════════════════════════════════════════════════════════════════════════
\section{Conclusions}
\label{sec:conclusions}
% ═════════════════════════════════════════════════════════════════════════════

We have presented a systematic comparison of five surrogate methods for
simulating two-dimensional tumour growth with spatial hypoxia under
fractionated radiotherapy.
Our main findings are:

\begin{enumerate}[leftmargin=2em]
  \item \textbf{Physics-informed methods outperform unconstrained black-box
        approaches} in generalisation beyond the training window.
        The NODE achieves the lowest in-window RMSE ($8.9\!\times\!10^{-3}$),
        while the PINN provides the most consistent accuracy across both
        windows ($\Delta\RMSE < 7\%$).

  \item \textbf{Simple MLP surrogates are unreliable for extrapolation.}
        The lack of physical regularisation causes error accumulation, with
        test RMSE $29\%$ higher than training RMSE in our experiment.
        Embedded physics (ODE baseline, PDE residual) is necessary for
        reliable long-horizon prediction.

  \item \textbf{The SuperNet teacher-student paradigm offers the best
        accuracy-efficiency trade-off} when training budget is limited.
        By distilling the PINN teacher into a fast scalar student, it achieves
        competitive accuracy at a fraction of the training cost (0.4\,s vs.\
        52.3\,s for the NODE).

  \item \textbf{Scalar surrogates provide orders-of-magnitude speedup at
        clinical grid resolutions} ($256^2$ or $64^3$), making them
        essential for ensemble forecasting, Bayesian inference, and
        treatment optimisation workflows.

  \item \textbf{Spatial averaging is a significant source of error
        during active irradiation.}
        The ODE baseline's high training RMSE ($1.44\!\times\!10^{-1}$) is
        attributable to heterogeneous hypoxia effects that spatial averaging
        cannot capture.
        Moment-based reduced models (beyond the zeroth spatial moment) are a
        promising avenue to address this limitation.
\end{enumerate}

The unified simulation framework released alongside this paper makes it
straightforward to add new engine implementations (by sub-classing
\texttt{BaseEngine}) and to plug them into the same visualisation, training,
and benchmarking infrastructure.

\printbibliography

\end{document}
